\documentclass{report}
\usepackage{amsmath,amssymb}
\setlength{\parindent}{0mm}
\setlength{\parskip}{1em}
\begin{document}
\begin{center}
\rule{6in}{1pt} \
%% SPEAKER NOW SANDERS:
%% {\large Christine Klymko \\
{\large Geoffrey Sanders \\
{\bf Detecting Highly-Cyclic Structure with Complex Eigenpairs}}

Lawrence Livermore National Laboratory \\ Center for Applied Scientific Computing \\ 7000 East Avenue \\ Livermore \\ Ca 94550
\\
% {\tt klymko1@llnl.gov}\\
{\tt sanders29@llnl.gov}\\
%Geoffrey Sanders\end{center}
Christine Klymko\end{center}

Highly 3- and 4-cyclic subgraph topologies are detectable via calculation
of eigenvectors associated with certain complex eigenvalues of Markov
propagators. We characterize this phenomenon theoretically to understand
the capabilities and limitations for utilizing eigenvectors in this
venture. We provide algorithms for approximating these eigenvectors and
give numerical results, both for software that utilizes complex
arithmetic and software that is limited to real arithmetic. Additionally,
we discuss the application of these techniques to motif detection.


\end{document}
